\section{Measurement problem and design}
\label{sec:design}

\subsection{The measurement chain}
A drawing feature is the last step of a chain.
The chain has a stylus, a digitiser, an operating system, a capture application, and the analysis software.
We call it the \emph{measurement chain}.
Two chains can report different feature values for the same movement.
The cause can be the sampling rate, the timestamps, the position scale, the filtering, or the noise.

We compare two chains.
Chain P is an iPad Pro with an Apple Pencil Pro.
Chain W is a Wacom Intuos Pro tablet with its pen, connected to a computer.
The same analysis software computes the features on both chains.
For brevity, we call the two chains device P and device W.

The question has one form for each feature.
Do device P and device W give equivalent values for the same movement, within a limit fixed in advance?
\emph{Equivalent} means that the difference between the devices lies inside the limit, with stated confidence, on every tested movement.
A test that finds no difference does not show equivalence.
We therefore use an equivalence test (Section~\ref{sec:stats}).

\subsection{Why a motion system and a camera}
A person cannot repeat a movement exactly.
A person also writes differently on glass than on a textured tablet surface~\cite{Gerth2016b,Alamargot2015}.
A comparison with people therefore mixes two effects: the difference between the chains and the change in the movement.
A motion system removes the second effect in part.
It executes the same \emph{program}, a prespecified list of timed moves, on both devices.

The \emph{commanded movement} is the movement that the program specifies.
The \emph{executed movement} is the movement that the stylus tip makes.
The two can differ, because the stylus holder bends under friction, and friction differs between the two surfaces.
A camera therefore records every run.
It measures the executed movement on both devices with one instrument, one length standard, and one clock.
For each run, we subtract the camera value of a feature from the device value.
This \emph{camera-referenced difference} removes the executed movement from the comparison.
What remains is the error of the chain, plus any error of the camera that differs between the two devices.
Section~\ref{sec:budget} treats this camera error.

\subsection{Design overview}
Fig.~\ref{fig:bench} shows the bench and the movement programs.
The motion system holds the stylus at a fixed angle and under a fixed load.
It executes 12 main programs: loops and spirals at two sizes and two speeds, one fast loops program, and three lines with a designed speed oscillation.
Three fault programs test the quality checks of the analysis software.

A \emph{run} is one execution of one program on one device.
A \emph{mounting} is one placement of a device in the fixture.
A \emph{block} is the set of all runs of one mounting.
Each device has six mountings, and each mounting has five runs of each main program.
The design therefore has 720 main runs and 108 fault runs.
The mounting is the unit of replication, because a new placement changes the alignment, the position on the active area, and the calibration.

Six features carry the equivalence decisions: motion duration, path length, mean speed, loop height, speed coefficient of variation (speed CV), and band power.
The band power is the power of the speed signal between 4 and 12~Hz, the band of most tremors.
The camera provides a reference for the first four features.
Each feature has an equivalence limit.
Three limits follow from a stated consequence for a screening model, and three are stated judgements (Section~\ref{sec:limits}).

All procedures, programs, limits, and analysis code are frozen before the main runs.
We call this step the \emph{lock}.
Nobody compares feature values between the devices until the last block is complete.
Two further steps use the bench results.
A correction maps the values of device W to the values of device P.
A consequence experiment shows what the measured difference does to a screening model.
It uses generated cohorts only.
The study has no participants and uses no data from people.
